% =============================================================================
% ForensiQ — Template de Architecture Decision Record (ADR)
% UC 21184 · Projeto de Engenharia Informática · UAb
%
% USO: Copiar este ficheiro, incrementar o número (ADR-XXXX),
%      preencher todos os campos e actualizar o estado.
% =============================================================================
\documentclass[a4paper,11pt]{article}

\usepackage[utf8]{inputenc}
\usepackage[T1]{fontenc}
\usepackage[portuguese]{babel}
\usepackage[top=2.5cm, bottom=2.5cm, left=3cm, right=2.5cm]{geometry}
\usepackage{booktabs}
\usepackage{longtable}
\usepackage{array}
\usepackage{parskip}
\usepackage{hyperref}
\usepackage{enumitem}

\hypersetup{
    colorlinks=true,
    linkcolor=black,
    urlcolor=blue,
    pdfauthor={João Rodrigues},
    pdftitle={ForensiQ — ADR-001 Template}
}

\newcolumntype{L}[1]{>{\raggedright\arraybackslash}p{#1}}

% =============================================================================
\begin{document}
% =============================================================================

\begin{center}
    {\LARGE\bfseries ADR-001 --- [Título curto da decisão]}\\[1.2em]
    \begin{tabular}{ll}
        \textbf{Data:}       & [Data] \\
        \textbf{Estado:}     & Proposto \textbar\ Aceite \textbar\ Substituído por ADR-XXX \\
        \textbf{Decisores:}  & [Nome(s)] \\
    \end{tabular}
\end{center}

\vspace{0.5em}
\hrule
\vspace{1.5em}

% =============================================================================
\section{Contexto}
% =============================================================================

% Qual é o problema ou situação que requer uma decisão?
% Que forças estão em jogo? (técnicas, de negócio, de equipa)

[Descrever o contexto que levou a esta decisão]

% =============================================================================
\section{Decisão}
% =============================================================================

% O que foi decidido, de forma afirmativa e directa.
% "Decidimos usar X" — não "Poderemos considerar X".

[Descrever a decisão tomada]

% =============================================================================
\section{Alternativas consideradas}
% =============================================================================

% Que outras opções foram avaliadas e porquê foram rejeitadas?

\begin{longtable}{L{5cm} L{9cm}}
    \toprule
    \textbf{Alternativa} & \textbf{Razão de rejeição} \\
    \midrule
    \endhead
    [Alternativa A] & [Porquê não] \\[0.4em]
    [Alternativa B] & [Porquê não] \\
    \bottomrule
\end{longtable}

% =============================================================================
\section{Consequências}
% =============================================================================

% O que muda com esta decisão?
% Que novas obrigações cria? Que problemas resolve? Que problemas pode criar?

\textbf{Positivas:}
\begin{itemize}[nosep, leftmargin=*]
    \item [Consequência positiva]
\end{itemize}

\textbf{Negativas / trade-offs:}
\begin{itemize}[nosep, leftmargin=*]
    \item [Trade-off aceite]
\end{itemize}

\vspace{2em}
\hrule
\vspace{0.5em}
{\small Para criar um novo ADR: copiar este ficheiro, incrementar o número,
preencher e actualizar o estado.}

\end{document}
